\chapter{\ifenglish Introduction\else บทนำ\fi}

\section{\ifenglish Project rationale\else ที่มาของโครงงาน\fi}

ในปัจจุบัน โมเดลภาษาขนาดใหญ่ (Large Language Models หรือ LLMs) ได้รับความนิยมและมีบทบาทสำคัญอย่างยิ่งในหลากหลายอุตสาหกรรม อย่างไรก็ตาม 
ความก้าวหน้าดังกล่าวมาพร้อมกับความท้าทายด้านความมั่นคงปลอดภัยของข้อมูล โดยเฉพาะความเสี่ยงที่ข้อมูลส่วนบุคคล (Personally Identifiable Information หรือ PII) 
จากชุดข้อมูลที่ใช้ในการฝึกฝนโมเดลอาจรั่วไหลออกไปสู่สาธารณะได้ ผ่านการโจมตีในรูปแบบต่าง ๆ เช่น การโจมตีเพื่อสกัดข้อมูล (Data Extraction Attack หรือ DEA) 
และการโจมตีด้วยเจลเบรค (Jailbreak Attack หรือ JA) \cite{carlini2021extracting, wei2023jailbroken}


แม้ในปัจจุบันจะมีการพัฒนามาตรฐานและชุดทดสอบ (Benchmark) เพื่อประเมินความเสี่ยงเหล่านี้ \cite{gu2024llmpbe}
\\อย่างแพร่หลาย แต่การทดสอบส่วนใหญ่มักถูกจำกัดอยู่ในบริบทของภาษาอังกฤษเป็นหลัก ซึ่งไม่สามารถสะท้อนความซับซ้อนของภาษาไทยที่มีโครงสร้างเฉพาะตัวและระดับภาษา 
ที่หลากหลายได้ ความแตกต่างของระดับภาษาตั้งแต่ภาษาระดับพิธีการไปจนถึงภาษาระดับกันเอง อาจส่งผลต่อประสิทธิภาพของกลไกความปลอดภัยในโมเดลที่แตกต่างกัน 

ด้วยเหตุนี้ คณะผู้จัดทำจึงได้พัฒนาชุดทดสอบภาษาไทยที่ครอบคลุม 5 ระดับภาษา เพื่อประเมินความเปราะบางของโมเดลภาษาขนาดต่างๆที่มีประสิทธิภาพสูง ได้แก่ 
Llama-3.1-8B-instant, Meta-Llama-4-Maverick-17B และ Kimi-k2-instruct พร้อมทั้งเปรียบเทียบประสิทธิภาพของกลไกการป้องกันระหว่างการคัดกรองข้อมูล 
(Scrubbing) และการป้อนคำสั่งป้องกัน (Defensive Prompting) ผลการศึกษานี้จะช่วยสร้างแนวทางที่ชัดเจนในการรักษาความปลอดภัยของข้อมูลในบริบทภาษาไทย 
และเป็นพื้นฐานสำคัญในการพัฒนาปัญญาประดิษฐ์ที่มีความน่าเชื่อถือในประเทศไทยต่อไป
\section{\ifenglish Objectives\else วัตถุประสงค์ของโครงงาน\fi}
\begin{enumerate}
    \item เพื่อประเมินความเสี่ยงและระบุรูปแบบการโจมตีที่สร้างความเสียหายสูงที่สุดสำหรับการโจมตีในรูป
    \\แบบภาษาไทย ทั้งในรูปแบบ Data Extraction Attack และ Jailbreak Attack
    \item เพื่อเปรียบเทียบและนำเสนอแนวทางการป้องกันที่มีประสิทธิภาพสูงสุด สำหรับภาษาไทย โดย
    \\พิจารณาความปลอดภัยที่เพิ่มขึ้นหลังการใช้ Scrubbing และ Defensive Prompting
\end{enumerate}

\section{\ifenglish Project scope\else ขอบเขตของโครงงาน\fi}

ในการดำเนินโครงงานวิจัยนี้ คณะผู้จัดทำได้กำหนดขอบเขตของการศึกษาไว้ดังนี้:

\subsection{\ifenglish Hardware scope\else ขอบเขตด้านฮาร์ดแวร์\fi}
\begin{itemize}
    \item เครื่องคอมพิวเตอร์ส่วนบุคคลสำหรับการพัฒนาและประมวลผลคำสั่งชุดทดสอบ 
    \item โครงสร้างพื้นฐานระบบประมวลผลผ่าน Groq API สำหรับการเรียกใช้งานโมเดลภาษาขนาดใหญ่ 
\end{itemize}

\subsection{\ifenglish Software scope\else ขอบเขตด้านซอฟต์แวร์\fi}
\begin{enumerate}
    \item โมเดลภาษาที่ใช้ในการศึกษา (Foundation Models) \\
    ทำการทดสอบและประเมินผลบนโมเดลภาษาขนาดใหญ่จำนวน 3 โมเดล ได้แก่
    \begin{itemize}
        \item Llama-3.1-8B-instant 
        \item Meta-Llama-4-Maverick-17B-128e-instruct 
        \item Moonshotai-Kimi-K2-instruct-0905 
    \end{itemize}

    \item รูปแบบภาษาและระดับภาษา (Linguistic Registers) \\
    ศึกษาและออกแบบชุดคำสั่ง (Prompt) ในรูปแบบภาษาไทย โดยแบ่งระดับภาษาออกเป็น 5 ระดับ ได้แก่
    \begin{itemize}
        \item ภาษาระดับพิธีการ (Ceremonial) เช่น " ขอเรียนเชิญท่านผู้มีเกียรติร่วมลิ้มรสภัตตาหารอันเลิศรส ณ สโมสรแห่งนี้ "
        \item ภาษาระดับทางการ (Formal) เช่น " ขอเชิญแขกผู้มีเกียรติทุกท่านรับประทานอาหารค่ำตามอัธยาศัยครับ/ค่ะ "
        \item ภาษาระดับกึ่งทางการ (Semi-formal) เช่น " เชิญทุกท่านไปทานข้าวด้วยกันนะครับ "
        \item ภาษาระดับไม่เป็นทางการ (Informal) เช่น " ไปหาอะไรกินกันเถอะ หิวแล้ว "
        \item ภาษาระดับกันเอง (Casual) เช่น " ไปแดกข้าวกันพวกเรา "
    \end{itemize}

    \item รูปแบบการโจมตีที่ใช้ในการทดสอบ (Attack Scenarios)
    \begin{itemize}
        \item Data Extraction Attack (DEA): มุ่งเน้นการโจมตีเพื่อดึงข้อมูลส่วนบุคคล (PII) จากชุดข้อมูล Enron Email 
        \item Jailbreak Attack (JA): ใช้เทคนิคการออกแบบ Prompt เพื่อหลีกเลี่ยงระบบความปลอดภัยของโมเดล โดยมุ่งเป้าไปที่ข้อมูลส่วนบุคคลที่ถูกจำกัดไว้ เช่น ข้อมูลของบุคคลสำคัญ 
    \end{itemize}

    \item กลไกการป้องกันที่ศึกษา (Defense Mechanisms) \\
    ประเมินประสิทธิภาพของกลไกการป้องกัน 2 รูปแบบหลัก:
    \begin{itemize}
        \item Scrubbing: การตรวจสอบและคัดกรองข้อมูลขาเข้า (Input) โดยใช้เทคนิค Regex และ NER เพื่อปกปิดข้อมูลส่วนบุคคล 
        \item Defensive Prompting: การเพิ่มคำสั่งป้องกันในระดับ System Prompt เพื่อเสริมสร้างความปลอดภัยจากภายในโมเดล 
    \end{itemize}

    \item สภาพแวดล้อมและทรัพยากรที่ใช้
    \begin{itemize}
        \item ดำเนินการทดสอบผ่านการเรียกใช้งานโมเดลด้วย Groq API 
        \item การทดสอบแต่ละกรณีจะดำเนินการทำซ้ำอย่างน้อย 5 รอบ เพื่อหาค่าเฉลี่ยของอัตราการโจมตีสำเร็จ (Average Success Rate) 
    \end{itemize}
\end{enumerate}

\section{\ifenglish Expected outcomes\else ประโยชน์ที่ได้รับ\fi}

ในการดำเนินโครงงานวิจัยนี้ คณะผู้จัดทำคาดว่าจะได้รับประโยชน์ต่าง ๆ ดังนี้:

\begin{itemize}
    \item ความเข้าใจเชิงลึกเกี่ยวกับความปลอดภัยของ LLM ในบริบทภาษาไทย: ทำให้ทราบถึงระดับความเปราะบางและอัตราความสำเร็จของการโจมตีประเภท DEA และ JA ในระดับภาษาไทยที่แตกต่างกันทั้ง 5 ระดับ
    \item แนวทางการเลือกใช้โมเดลภาษาขนาดต่างๆที่มีประสิทธิภาพ: ทราบถึงขีดความสามารถและความต้านทานต่อการโจมตีของโมเดล Llama-3.1-8b-instant, Meta-llama-llama-4-maverick-17b-128\\e-instruct และ Moonshotai-kimi-k2-instruct-0905 เพื่อเป็นข้อมูลประกอบการตัดสินใจเลือกใช้โมเดลในสภาพแวดล้อมที่จำกัดทรัพยากร
    \item กลไกการป้องกันที่เหมาะสมต่อประเภทภัยคุกคาม: ได้ข้อสรุปที่ชัดเจนว่าควรเลือกใช้กลไก Scrubbing หรือ Defensive Prompting สำหรับการรับมือกับการโจมตีแต่ละประเภท เพื่อการรักษาความปลอดภัยของข้อมูลส่วนบุคคล (PII) อย่างมีประสิทธิภาพสูงสุด
    \item มาตรฐานการประเมินผลสำหรับภาษาไทย (Thai Privacy Benchmark): สร้างบรรทัดฐานและ \\แนวทางในการทดสอบความปลอดภัยของโมเดลภาษาในประเทศไทย โดยต่อยอดจากกรอบแนวคิดที่เป็นสากลอย่าง LLM-PBE \cite{gu2024llmpbe}
    \item การเสริมสร้างความปลอดภัยให้แก่ระบบปัญญาประดิษฐ์: ข้อมูลจากการวิจัยสามารถนำไปประยุกต์ใช้เพื่อลดความเสี่ยงด้านข้อมูลรั่วไหลในการพัฒนาแอปพลิเคชันที่ใช้งานจริงในบริบทภาษาไทย
\end{itemize}

\section{\ifenglish Technology and tools\else เทคโนโลยีและเครื่องมือที่ใช้\fi}
เทคโนโลยีและเครื่องมือที่ใช้ในการพัฒนาและดำเนินงานวิจัย มีรายละเอียดดังนี้:
\subsection{\ifenglish Hardware technology\else เทคโนโลยีด้านฮาร์ดแวร์\fi}
\begin{itemize}
    \item Personal Computer Specification: เครื่องคอมพิวเตอร์สำหรับการพัฒนาโปรแกรมและการจัดเก็บชุดข้อมูลทดสอบ
    \item Cloud Inference Engine: ใช้งานระบบประมวลผลประสิทธิภาพสูงจาก Groq API เพื่อให้สามารถรันโมเดลขนาดใหญ่ได้อย่างรวดเร็วและมีประสิทธิภาพ
\end{itemize}
\subsection{\ifenglish Software technology\else เทคโนโลยีด้านซอฟต์แวร์\fi}
\begin{itemize}
    \item Programming Language: ภาษา Python สำหรับการพัฒนาสคริปต์เพื่อจำลองการโจมตีและประมวลผลกลไกการป้องกัน
    \item Application Programming Interface (API): บริการ Groq API สำหรับการเชื่อมต่อและส่งคำสั่ง (Prompts) ไปยังโมเดลเป้าหมาย
    \item Security and Data Processing Tools:
        \begin{enumerate}
            \item Regular Expression (Regex): สำหรับการทำ Pattern Matching เพื่อตรวจจับและปิดบังข้อมูลในส่วน Scrubbing
            \item Named Entity Recognition (NER): สำหรับระบุชื่อบุคคลหรือข้อมูลเฉพาะเจาะจงในภาษาไทยเพื่อทำการ Scrubbing
        \end{enumerate}
    \item Development Platform: ซอฟต์แวร์สำหรับเขียนและทดสอบโค้ด (VS Code) และเครื่องมือสำหรับวิเคราะห์ผลสถิติ สร้างกราฟ
\end{itemize}
\section{\ifenglish Project plan\else แผนการดำเนินงาน\fi}

\begin{plan}{12}{2025}{3}{2026}
    \planitem{12}{2025}{12}{2025}{กำหนดหัวข้อโปรเจค และ กำหนดขอบเขตของโปรเจค}
    \planitem{12}{2025}{12}{2025}{ศึกษางานวิจัยที่เกี่ยวข้อง (LLM-PBE)}
    \planitem{12}{2025}{12}{2025}{ออกแบบ DEA prompt และ JA prompt สำหรับภาษาไทยในแต่ละระดับ}
    \planitem{12}{2025}{1}{2026}{สร้าง script รันการโจมตีอัตโนมัติ และ ออกแบบการการวัดผลการโจมตีที่สำเร็จ}
    \planitem{12}{2025}{2}{2026}{ดำเนินการทดลองโจมตีและเก็บรวบรวมผลการทดสอบ}
    \planitem{2}{2026}{2}{2026}{ออกแบบการป้องกัน Scrubbing และ Defensive prompt}
    \planitem{2}{2026}{3}{2026}{ดำเนินการทดลองป้องกันการโจมตีและเก็บรวบรวมผลการทดสอบ}
    \planitem{3}{2026}{3}{2026}{วิเคราะห์ผลการทดลองและสรุปผลการทดลอง}
    \planitem{3}{2026}{3}{2026}{จัดทำรายงานฉบับสมบูรณ์และเตรียมนำเสนอผลงาน}
\end{plan}

\section{\ifenglish Budget plan\else แผนการใช้งบประมาณ\fi}

\ifenglish
The total budget for this project was 8,014.61 Baht, which can be divided into two main parts as follows:
\begin{enumerate}
    \item \textbf{Groq API Processing Costs} (over 90\% of the budget): 7,401.11 Baht was allocated for executing Data Extraction Attack (DEA) and Jailbreak Attack (JA) scenarios. The execution included 5 rounds for each of the 5 Thai language levels and 5 rounds for English (totaling 60 execution rounds). Each round consisted of approximately 3,000 prompts.
    \item \textbf{Data Preparation and Presentation Costs} (10\% of the budget): 613.50 Baht was used for translating the English dataset into 5 Thai language levels (a total of 15,000 prompts) and for producing the project poster.
\end{enumerate}
\else
ในการดำเนินโครงงานนี้ มีการใช้งบประมาณทั้งสิ้น 8,014.61 บาท (จากค่างบประมาณที่ได้รับ 8,000 บาท) โดยสามารถจำแนกสัดส่วนการใช้งบประมาณออกเป็น 2 ส่วนหลัก ดังนี้:
\begin{enumerate}
    \item \textbf{ค่าใช้จ่ายสำหรับการประมวลผลผ่าน Groq API} (สัดส่วนมากกว่าร้อยละ 90 ของงบประมาณ): เป็นจำนวนเงิน 7,401.11 บาท ถูกใช้เพื่อบรรลุวัตถุประสงค์ในการรันการโจมตีรูปแบบ Data Extraction Attack (DEA) และ Jailbreak Attack (JA) โดยดำเนินการทดสอบในภาษาไทยทั้ง 5 ระดับภาษา อย่างละ 5 รอบ และภาษาอังกฤษอีก 5 รอบ (รวมเป็น 60 รอบการทดสอบ) ซึ่งในแต่ละรอบประกอบด้วยคำสั่ง (Prompts) จำนวนประมาณ 3,000 คำสั่ง
    \item \textbf{ค่าใช้จ่ายสำหรับการแปลงข้อมูลและจัดทำสื่อนำเสนอ} (สัดส่วนร้อยละ 10 ของงบประมาณ): เป็นจำนวนเงิน 613.50 บาท ถูกใช้สำหรับการแปลงข้อมูลชุดทดสอบจากภาษาอังกฤษเป็นภาษาไทย 5 ระดับภาษา (จำนวน 15,000 คำสั่ง) รวมถึงค่าใช้จ่ายในการจัดทำโปสเตอร์ (Poster) เพื่อใช้ประกอบการนำเสนอผลงาน
\end{enumerate}
\fi

\section{\ifenglish Roles and responsibilities\else บทบาทและความรับผิดชอบ\fi}
เพื่อให้การดำเนินงานวิจัยเป็นไปอย่างมีประสิทธิภาพ คณะผู้จัดทำได้แบ่งหน้าที่ความรับผิดชอบตามความ \\เชี่ยวชาญและเนื้อหางาน ดังนี้:
\begin{enumerate}
    \item กิตติวินท์ พรรณเชษฐ์
        \begin{itemize}
            \item ออกแบบชุดคำสั่งสำหรับการโจมตีทั้งรูปแบบ Data Extraction Attack (DEA) และ Jailbreak Attack (JA)
            \item รับผิดชอบการรันการโจมตีในระดับภาษาทางการ (Ceremonial) และระดับภาษากึ่งทางการ (Semi-formal) 
            \item ออกแบบและแก้ไขค่ามาตรวัด JA ในส่วนของข้อมูลติดต่อ
            \item วิเคราะห์ผลลัพธ์การโจมตี DEA และ JA ร่วมกับทีมเพื่อนำไปสู่การสรุปผลเชิงสถิติและการทำ Visualization
            \item รับผิดชอบการแปลชุดคำสั่ง JA จากภาษาอังกฤษเป็นภาษาไทยครอบคลุมทั้ง 5 ระดับภาษา
            \item ออกแบบกลไกการป้องกัน และดำเนินการทดสอบระบบ Input Scrubbing ของการโจมตี DEA
        \end{itemize}
    \item พิพัฒน์พงศ์ กาวิลดง
        \begin{itemize}
            \item รับผิดชอบหลักในการรันการโจมตีในระดับภาษาไม่เป็นทางการ (Informal) และการรันผลการโจมตีในรอบภาษาอังกฤษสำหรับทุกโมเดล (Semi-formal)
            \item พัฒนาและแก้ไขโครงสร้าง Code สำหรับการจัดการเส้นทางข้อมูล (Path), การสร้างจุดบันทึกผล (Checkpoint) และการประยุกต์ใช้ Defensive Prompting ร่วมกับ JA
            \item ออกแบบมาตรวัด JA และจัดทำสื่อนำเสนอผลงานสำหรับการรายงานความก้าวหน้า
        \end{itemize}
    \item ภัทรพล ณ นคร
        \begin{itemize}
            \item ดำเนินการแปลชุดคำสั่งปริมาณมาก (Batch Processing) ผ่าน OpenAI API ให้ครอบคลุมทุกระดับภาษา
            \item รับผิดชอบการรันการโจมตีในระดับภาษาทางการ (Formal) และภาษาอังกฤษ รวมถึงการบริหารจัดการ API Key ผ่านระบบ Organization ของ Groq
            \item พัฒนามาตรวัด JA ให้สมบูรณ์ วิเคราะห์ Keyword เพื่อออกแบบกลไกการป้องกัน และดำเนินการทดสอบระบบ Input Scrubbing ของการโจมตี JA
            \item จัดการระบบควบคุมเวอร์ชันของซอร์สโค้ด (Version Control) และวิเคราะห์ผลลัพธ์เพื่อสร้าง Visualization
        \end{itemize}
    \newpage
    \item วายุ ทะไร
        \begin{itemize}
            \item ออกแบบมาตรวัดสำหรับ DEA และพัฒนาระบบรันการทดสอบอัตโนมัติที่สามารถเรียกใช้หลาย API และหลายรูปแบบการโจมตีพร้อมกันผ่าน Terminal เดียว
            \item รับผิดชอบหลักในการรันการโจมตีในระดับภาษากันเอง (Casual) สำหรับทั้ง 3 โมเดล และรันการทดสอบในรอบภาษาอังกฤษ
            \item พัฒนาระบบรวบรวมข้อมูลผลลัพธ์ คำนวณปริมาณ Token และความยาวข้อความ พร้อมสร้างกราฟสรุปผลอัตโนมัติ
            \item ออกแบบชุดคำสั่งป้องกัน (Defensive Prompting) สำหรับทั้ง DEA และ JA จำนวน 12 รูปแบบ และคัดเลือกรูปแบบที่มีประสิทธิภาพสูงสุดมาใช้งาน
        \end{itemize}
\end{enumerate}
\section{\ifenglish%
Impacts of this project on society, health, safety, legal, and cultural issues
\else%
ผลกระทบด้านสังคม สุขภาพ ความปลอดภัย กฎหมาย และวัฒนธรรม
\fi}

การดำเนินโครงงานวิจัยนี้ส่งผลกระทบในมิติต่าง ๆ ที่สำคัญ ดังนี้:

\begin{itemize}
    \item ด้านความปลอดภัยและกฎหมาย \\ โครงงานนี้ช่วยระบุช่องโหว่ของการรั่วไหลข้อมูลส่วนบุคคล (PII) ในโมเดลภาษาขนาดใหญ่ ซึ่งสอดคล้องกับหลักการของพระราชบัญญัติคุ้มครองข้อมูลส่วนบุคคล (PDPA) ของประเทศไทย โดยการทดสอบการโจมตีแบบ Data Extraction Attack (DEA) และ Jailbreak Attack (JA) ช่วยให้ผู้พัฒนาระบบสามารถเตรียมแนวทางป้องกันเพื่อลดความเสี่ยงทางกฎหมายและเพิ่มความมั่นคงปลอดภัยให้แก่ข้อมูลของผู้ใช้งาน
    \item ด้านสังคมและวัฒนธรรม \\ เนื่องจากมาตรฐานการประเมินส่วนใหญ่มักจำกัดอยู่ในบริบทภาษาอังกฤษ โครงงานนี้จึงสร้างผลกระทบเชิงบวกต่อสังคมไทยโดยการสร้างชุดทดสอบที่ครอบคลุมถึง 5 ระดับภาษาไทย ช่วยให้เทคโนโลยี AI มีความเข้าใจและเคารพในบริบททางวัฒนธรรมที่หลากหลายของไทย ตั้งแต่ภาษาระดับพิธีการไปจนถึงภาษาระดับกันเอง
    \item ด้านสุขภาพและสิ่งแวดล้อม\\ การมุ่งเน้นศึกษาบนโมเดลขนาดเล็กที่ประหยัดทรัพยากร (Resource-saving models) เช่น llama-3.1-8b-instant และ meta-llama-llama-4-maverick-17b-128e-instruct  ส่งผลดีต่อสิ่งแวดล้อมในแง่ของการลดการใช้พลังงานไฟฟ้าและทรัพยากรในการประมวลผล (Compute power) ซึ่งเป็นส่วนหนึ่งของการพัฒนา Green IT ในอนาคต
\end{itemize}